% ============================================================
%  第七章 总结与展望
% ============================================================
\chapter{总结与展望}

\section{工作总结}

本文设计并实现了一套基于Spring Boot的高校社团管理系统，主要完成了以下工作：

\begin{enumerate}
  \item \textbf{系统架构设计}：采用模块化单体架构，将业务逻辑按职责拆分为8个独立模块，通过清晰的模块边界实现低耦合、高内聚的代码组织；
  \item \textbf{安全体系建设}：基于Spring Security 6.x与JWT双令牌机制实现无状态认证，结合RBAC权限模型、登录限流、Token黑名单与AOP审计日志，构建了完善的安全防护体系；
  \item \textbf{实时交互能力}：通过WebSocket（STOMP/SockJS）实现活动签到实时广播与系统通知即时推送，显著提升了用户交互体验；
  \item \textbf{AI智能服务集成}：引入基于LangChain与Chroma的RAG知识库问答系统，以及融合SVD协同过滤与语义内容推荐的混合推荐算法，为用户提供智能化服务；
  \item \textbf{全面测试验证}：通过单元测试、接口测试、安全测试与性能测试，验证了系统功能的正确性、安全性与性能达标情况。
\end{enumerate}

系统的主要创新点在于：将AI智能服务（RAG问答与混合推荐）与传统社团管理系统深度融合，并通过模块化架构设计为未来的微服务化演进预留了扩展空间。

\section{不足与改进}

当前系统存在以下不足：

\begin{enumerate}
  \item \textbf{推荐算法冷启动问题}：新用户缺乏历史行为数据，SVD协同过滤效果有限，需引入基于用户画像的冷启动策略；
  \item \textbf{缺乏移动端支持}：系统目前仅支持PC浏览器访问，未提供移动端原生应用或响应式适配；
  \item \textbf{单点部署风险}：当前为单实例部署，缺乏高可用集群配置，存在单点故障风险；
  \item \textbf{RAG知识库更新延迟}：知识库同步依赖手动触发，无法实时反映最新数据变化。
\end{enumerate}

\section{未来展望}

基于当前系统的基础，未来可从以下方向进行改进与扩展：

\begin{enumerate}
  \item \textbf{微服务化演进}：随着用户规模增长，可将各业务模块拆分为独立微服务，引入Spring Cloud Gateway、Nacos等组件，提升系统的水平扩展能力；
  \item \textbf{AI功能深化}：优化推荐算法，引入强化学习动态调整推荐策略；升级RAG系统，支持多模态知识库（图片、视频）；
  \item \textbf{移动端适配}：开发基于uni-app或React Native的移动端应用，提升用户访问便捷性；
  \item \textbf{多租户支持}：扩展系统架构以支持多所高校同时使用，实现数据隔离与个性化配置；
  \item \textbf{自动化运维}：引入CI/CD流水线（GitHub Actions + Docker）与Kubernetes容器编排，实现自动化构建、测试与部署。
\end{enumerate}
